\documentclass[10pt,a4paper]{amsart}
\usepackage[margin=19mm]{geometry}
\usepackage{amsmath,amssymb,mathtools}
\usepackage[scr=rsfs,cal=euler]{mathalfa}
\usepackage{xfrac}
\usepackage[unicode,hidelinks,bookmarks=false]{hyperref}
\newcommand{\ie}{i.e.\ }
\newcommand{\cf}{cf.\ }
\newcommand{\resp}{respectively\ }
\newcommand{\Subsection}{\subsection}
\providecommand{\nobreakdash}{\nobreak\hbox{-}}
\setlength{\emergencystretch}{2em}
\multlinegap0pt
\usepackage{amssymb}
\usepackage{esint}
\usepackage{mathtools}
\usepackage[shortlabels]{enumitem}
\usepackage{tikz}
\newtheorem{theorem}{Theorem}
\newtheorem{lemma}{Lemma}
\newcommand{\bbbr}{\mathbb R}
\newcommand{\bbbz}{\mathbb Z}
\def\diam{\operatorname{diam}}
\title{The Trudinger inequality is true for $M^{1,s}$ spaces}
\author{Ryan Alvarado, Ahmed Dughayshim, Piotr Haj\l{}asz}
\date{Source: arXiv:2609.09101v1 (CC BY 4.0)}
\begin{document}
\maketitle

\noindent\emph{Source and licence.} The Trudinger inequality is true for $M^{1,s}$ spaces. Version arXiv:2609.09101v1 (CC BY 4.0), distributed by arXiv under the Creative Commons Attribution 4.0 International licence, http://creativecommons.org/licenses/by/4.0/, as recorded in the arXiv licence metadata for this exact version and shown on the abstract page https://arxiv.org/abs/2609.09101v1. Full licence notice: \texttt{licenses/arxiv-2609.09101-CC-BY-4.0.txt}.\par\smallskip

\noindent\textit{Editorial scope.} Counted unit: the article's theorem T2 (source0.tex lines 255--319). The article's complete original proof of it is delivered with it (source0.tex lines 845--1375, one \texttt{proof} environment of 531 lines, which the article places away from the statement). No mathematics is written, repaired or re-derived: the statement and the proof are the article's own lines, in the article's own notation, and no retained line is substituted or re-flowed.  Retained uncounted as the source-local context the proof needs: lines 187--232; lines 519--525; lines 535--538; lines 540--545; lines 562--567; lines 584--588; lines 590--595; lines 660--669; lines 701--709; lines 748--752; lines 781--789.  Nothing was omitted from any retained span, so the package carries no omission record.  No retained line contains a citation, so the wrapper carries no bibliography and no bibliography entry.  No retained line contains a figure environment, an image-inclusion command or drawing code, so no image is delivered and the package declares no asset dependency.  Editorial formatting only, in addition: an AMS \texttt{amsart} preamble that restates the article's own declarations and macros used by the retained lines, namely the environments \texttt{theorem}, \texttt{lemma} on separate counters, together with the standard packages those lines need.  The retained environments print in this document as Theorem 1, Lemma 1 in order of appearance, whereas the article prints the same items with the numbering of its own full document.  The complete native source of the article as distributed in the arXiv e-print for this exact version is bundled unchanged as \texttt{sources/209777/source0.tex}.
\begin{theorem}
\label{T1}
Let $(X,d)$ be bounded, $0<\mu(X)<\infty$, and assume that for some $s>0$ and $b>0$,
\begin{equation}
\label{eq1}
 \mu(B(x,r))\ge br^s,
 \qquad x\in X,\quad 0<r\leq\diam X.
\end{equation}
Let $u\in M^{1,p}(X)$ and let $g\in D(u)\cap L^p(X)$.
\begin{enumerate}
\item If $0<p<s$ and $p^*=\frac{sp}{s-p}$,
then
\begin{equation}
\label{eq57}
 \inf_{\gamma\in\mathbb R}\Vert u-\gamma\Vert_{p^*}
 \leq C_{s,p}b^{-1/s}\Vert g\Vert_p.
\end{equation}
If $p^*\ge1$, then $u_X$ is defined and one may take $\gamma=u_X$, after changing the constant to $2C_{s,p}$. If $1\leq p<s$, one can take in \eqref{eq57}
\begin{equation}
\label{eq16}
C_{s,p}=64\cdot 2^s\left(\frac{s-1}{s-p}\right)^{1/s'},
\end{equation}
\item If $p=s>1$, then
\begin{equation}
\label{eq17}
\fint_X
 \exp\!\left[
 c_s\left(
 \frac{b^{1/s}|u-u_X|}{\Vert g\Vert_{s}}
 \right)^{\frac{s}{s-1}}
 \right]d\mu
 \leq C_s.
\end{equation}
If $p=s$ and $0<s\leq1$, then
\[
 \Vert u-u_X\Vert_{\infty}
 \leq C_s b^{-1/s}\Vert g\Vert_{s}.
\]

\item If $s<p<\infty$, then
\[
 \Vert u-u_X\Vert_{\infty}
 \leq C_{s,p}b^{-1/s}\mu(X)^{1/s-1/p}\Vert g\Vert_p.
\]
\end{enumerate}
\end{theorem}

\begin{equation}
\label{eq9}
a_k\leq a_{k_o}
+2^{1+1/s}b^{-1/s}\mu(X)^{-1/p}\Vert g\Vert_p\sum_{i=k_o+1}^k 2^im_{i-1}^{1/s},
\qquad
k> k_o.
\end{equation}

\begin{equation}
\label{eq6}
 r_k=\left(\frac{2m_{k-1}}b\right)^{1/s}.
\end{equation}

\begin{equation}
\label{eq7}
d(x,y)\leq r_k
\qquad
\text{for some } y\in E_{k-1}.
\end{equation}

\begin{equation}
\label{eq8}
\sum_{k\in\mathbb Z}
2^{kp}m_{k-1}
\leq \frac{2^p}{1-2^{-p}} \mu(X).
\end{equation}

\begin{equation}
\label{eq10}
a_k\leq \Gamma_{s,p}\,b^{-1/s}\mu(X)^{-1/p^*}\Vert g\Vert_p\,2^{kp/p^*},
\qquad k\ge k_o,
\end{equation}

\begin{equation}
\label{eq15}
\Gamma_{s,p}\leq 32\left(\frac{s-1}{s-p}\right)^{1/s'}
\qquad
\text{when } 1\leq p<s.
\end{equation}

\begin{equation}
\label{eq39}
\sum_{i=k_o+1}^k2^im_{i-1}^{1/s}
\leq
\left(\sum_{i=k_o+1}^k2^{is}m_{i-1}\right)^{1/s}
\leq
2^{kp/p^*}
\left(\sum_{i=k_o+1}^k2^{ip}m_{i-1}\right)^{1/s}
\leq C_{s,p}2^{kp/p^*}\mu(X)^{1/s}.
\end{equation}

\begin{equation}
\label{eq11}
 \sum_{k\geq k_o}2^{kp}\mu(F_k)
 \leq
2^{s+1}\mu(X)+ 
2^p\mu(X)\Vert g\Vert_p^{-p}
\sum_{k>k_o}\int_{F_k}g^p\,d\mu
 \leq 2^{s+2}\mu(X).
\end{equation}

\begin{equation}
\label{eq12}
 a_k\leq C_s b^{-1/s}\Vert g\Vert_{s} k^{1/s'},
 \qquad k\geq k_o.
\end{equation}

\begin{lemma}
\label{T6}
If $q>1$ and $C>0$, then
\[
\fint_X\exp\left(\frac{C}{2^q}|u-u_X|^q\right)\, d\mu\leq
\inf_{\gamma\in\bbbr}
\fint_X\exp(C|u-\gamma|^q)\, d\mu.
\]
\end{lemma}

\begin{theorem}
\label{T2}
Let $0<p<\infty$, $1<\sigma\leq2$, and let $B=B(z,R)$ be any ball of radius $R$. Assume that for some $s,b,R_o>0$,
\begin{equation}
\label{eq20}
 \mu(B(x,r))\geq br^s,
 \qquad x\in\sigma B,\quad 0<r\leq R_o.
\end{equation}
Set
\[
 L=\min\{R_o,(\sigma-1)R\}.
\]
Let $u\in M^{1,p}(\sigma B)$ and let $g\in D(u)\cap L^p(\sigma B)$.
\begin{enumerate}
\item If $0<p<s$ and $p^*=sp/(s-p)$, then
\[
 \inf_{\gamma\in\bbbr}
 \left(\fint_B|u-\gamma|^{p^*}\,d\mu\right)^{1/p^*}
 \leq C_{s,p}\frac{R}{L}\,b^{-1/p}L^{1-s/p}
 \Vert g\Vert_{L^p(\sigma B)}.
\]
If $p^*\geq1$, then one may take $\gamma=u_B$, after changing the constant. If $1\leq p<s$, then, with $s'=s/(s-1)$, we get the sharper estimate
\begin{equation}
\label{eq54}
 \left(\fint_B|u-u_B|^{p^*}\,d\mu\right)^{1/p^*}
 \leq C_s\left(
 \frac{R}{L}+\left(\frac{s-1}{s-p}\right)^{1/s'}
 \right)b^{-1/p}L^{1-s/p}
 \Vert g\Vert_{L^p(\sigma B)}.
\end{equation}

\item If $p=s>1$, then
\[
 \fint_B
 \exp\!\left[
 c_s\left(
 \frac{L}{R}\frac{b^{1/s}|u-u_B|}
 {\Vert g\Vert_{L^s(\sigma B)}}
 \right)^{\frac{s}{s-1}}
 \right]d\mu
 \leq C_s.
\]
If $p=s\leq1$, then $u|_B$ has a  uniformly continuous representative that satisfies
\begin{equation}
\label{eq osc}
 \operatorname{osc}_B u
 \leq C_s\frac{R}{L}\,b^{-1/s}
 \Vert g\Vert_{L^s(\sigma B)}.
\end{equation}

\item If $p>s$, then $u|_B$ has a H\"older continuous representative of order $1-s/p$ that satisfies
\[
 \operatorname{osc}_B u
 \leq C_{s,p}\frac{R}{L}\,b^{-1/p}L^{1-s/p}
 \Vert g\Vert_{L^p(\sigma B)},
\]
and
\[
 |u(x)-u(y)|
 \leq C_{s,p}\frac{R}{L}\,b^{-1/p}d(x,y)^{1-s/p}
 \Vert g\Vert_{L^p(\sigma B)}
\]
for all $x,y\in B$.
\end{enumerate}
\end{theorem}

\begin{proof}[Proof of Theorem~\ref{T2}]
We may assume that $\Vert g\Vert_{L^p(\sigma B)}>0$, since otherwise $u$ is constant almost everywhere on $\sigma B$ and all the conclusions are immediate. Enlarging a null set $N$, if necessary, we may also assume that $g$ is finite on $\sigma B\setminus N$ and that
\[
 |u(x)-u(y)|\leq d(x,y)(g(x)+g(y))
\]
for every $x,y\in\sigma B\setminus N$.

We define the level sets exactly as in the proof of Theorem~\ref{T1}, with $\mu(B)$ in place of $\mu(X)$. Namely, for $k\in\bbbz$ let
\[
 E_k=\left\{x\in\sigma B\setminus N:\
 g(x)\leq 2^k\mu(B)^{-1/p}\Vert g\Vert_{L^p(\sigma B)}\right\},
 \qquad
 m_k=\mu(\sigma B\setminus E_k).
\]
The sets $E_k$ are increasing. Enlarging $N$ by a countable union of null sets, if necessary, we may assume that whenever $m_k=0$, we have $E_k=\sigma B\setminus N$, and hence $E_j=\sigma B\setminus N$ for every $j\geq k$. Since $g$ is finite on $\sigma B\setminus N$,
\[
 \sigma B\setminus N=\bigcup_{k\in\bbbz}E_k.
\]
Chebyshev's inequality gives
\begin{equation}
\label{eq21}
 m_k\leq 2^{-kp}\mu(B).
\end{equation}
Similar to the proof of Theorem~\ref{T1}, we shall need the following estimate:
\begin{equation}
\label{eq22}
 \sum_{k\in\bbbz}2^{kp}m_{k-1}
 \leq \frac{2^p}{1-2^{-p}}\mu(B).
\end{equation}
This is proven using the same Fubini calculation as in \eqref{eq8}, where the integration is over $\sigma B$ and we normalize by $\mu(B)$ instead of $\mu(X)$.


Recall that
\[
 L=\min\{R_o,(\sigma-1)R\}.
\]
Since $L\leq R$ and $L\leq R_o$, applying \eqref{eq20} at the center $z$ of $B$ gives
\begin{equation}
\label{eq23}
 \mu(B)\geq\mu(B(z,L))\geq bL^s.
\end{equation}
Put
\[
 \Lambda_{s,p}=\left(\frac{2^{1/s}}{1-2^{-p/s}}\right)^{s/p},
\]
and let $k_o$ be the smallest integer such that
\[
 \Lambda_{s,p}^{p/s}\,
 b^{-1/s}\mu(B)^{1/s}2^{-k_op/s}\leq L.
\]
Such an integer exists because the left hand side tends to zero as $k_o\to\infty$. By the minimality of $k_o$,
\begin{equation}
\label{eq24}
 \Lambda_{s,p}b^{-1/p}\mu(B)^{1/p}L^{-s/p}
 \leq2^{k_o}
 <2\Lambda_{s,p}b^{-1/p}\mu(B)^{1/p}L^{-s/p}.
\end{equation}
Notice that \eqref{eq23} and the left hand inequality in \eqref{eq24} imply $2^{k_o}>1$, so $k_o\geq1$. Hence, by \eqref{eq21},
\[
 m_{k_o}\leq2^{-k_op}\mu(B)<\mu(B),
\]
and therefore $E_{k_o}\cap B\ne\varnothing$. Fix $x_o\in E_{k_o}\cap B$ and put $\gamma=u(x_o)$. On $E_{k_o}$ the function $u$ is
$2^{k_o+1}\mu(B)^{-1/p}\Vert g\Vert_{L^p(\sigma B)}$-Lipschitz. Since $1<\sigma\leq2$,
$\diam(\sigma B)\leq4R$, and hence
\begin{equation}
\label{eq25}
 a_{k_o}:=\sup_{E_{k_o}}|u-\gamma|
 \leq2^{k_o+3}R\mu(B)^{-1/p}\Vert g\Vert_{L^p(\sigma B)}
 \leq 16\Lambda_{s,p}\frac RL\,b^{-1/p}L^{1-s/p}
 \Vert g\Vert_{L^p(\sigma B)}.
\end{equation}
In the last inequality we used the upper bound in \eqref{eq24}.

For $j>k_o$ define
\[
 r_j=\left(\frac{2m_{j-1}}b\right)^{1/s}.
\]
By \eqref{eq21},
\[
 r_j\leq2^{1/s}b^{-1/s}\mu(B)^{1/s}2^{-(j-1)p/s}.
\]
Therefore, by the definition of $k_o$,
\begin{equation}
\label{eq26}
 \sum_{j=k_o+1}^{\infty}r_j
 \leq\frac{2^{1/s}}{1-2^{-p/s}}\,
 b^{-1/s}\mu(B)^{1/s}2^{-k_op/s}
 \leq L.
\end{equation}
In particular, $r_j\leq L\leq R_o$ for every $j>k_o$.

For $x\in E_k\cap B$, $k>k_o$, we use the construction from
\eqref{eq6}--\eqref{eq7} to choose points $x_k=x$ and $x_j\in E_j$,
$k_o\leq j<k$, with $d(x_j,x_{j-1})\leq r_j$. However, this time we must also verify that every ball used in this construction lies inside $\sigma B$.
Indeed, suppose that $x_j,\ldots,x_k$ have already been chosen. 
If $m_{j-1}=0$, then
$E_{j-1}=E_j=\sigma B\setminus N$ and $r_j=0$, so taking
$x_{j-1}=x_j$ gives the desired step. We may therefore assume
that $m_{j-1}>0$. Then, for $w\in B(x_j,r_j)$, we use $x\in B$ and \eqref{eq26} to estimate,
\[
 d(w,z)<R+\sum_{i=j}^k r_i\leq R+L\leq\sigma R.
\]
Thus $B(x_j,r_j)\subset\sigma B$. Since $0<r_j\leq R_o$, \eqref{eq20} gives
\[
 \mu(B(x_j,r_j))\geq br_j^s=2m_{j-1}>m_{j-1}.
\]
It follows that $B(x_j,r_j)\cap E_{j-1}\neq\varnothing$, and we may therefore choose $x_{j-1}$ in this intersection. This completes the construction.


Since $E_{j-1}\subset E_j$, both $x_j$ and $x_{j-1}$ belong to $E_j$. Hence
\begin{align*}
|u(x_j)-u(x_{j-1})|
&\leq2^{j+1}\mu(B)^{-1/p}\Vert g\Vert_{L^p(\sigma B)}r_j
\\
&=2^{1+1/s}b^{-1/s}\mu(B)^{-1/p}\Vert g\Vert_{L^p(\sigma B)}
2^jm_{j-1}^{1/s}.
\end{align*}
Combining this estimate with the definition of $a_{k_o}$, we obtain the basic estimate
\begin{equation}
\label{eq27}
 \sup_{E_k\cap B}|u-\gamma|
 \leq a_{k_o}
 +2^{1+1/s}b^{-1/s}\mu(B)^{-1/p}\Vert g\Vert_{L^p(\sigma B)}
 \sum_{j=k_o+1}^k2^jm_{j-1}^{1/s},
 \qquad k>k_o.
\end{equation}
This is the local counterpart of the estimate \eqref{eq9}. 
Since \eqref{eq22} provides the same summability
bound as \eqref{eq8}, with $\mu(B)$ in place of $\mu(X)$,
we will refer to the corresponding arguments
in the proof of Theorem~\ref{T1} when estimating the sum
on the right-hand side of \eqref{eq27}.


\medskip
\noindent
{\sc The subcritical case $0<p<s$.}
Let $p^*=sp/(s-p)$. Our first aim is to prove that
\begin{equation}
\label{eq28}
 \sup_{E_k\cap B}|u-\gamma|
 \leq C_{s,p}\frac RL\,b^{-1/s}\mu(B)^{-1/p^*}
 \Vert g\Vert_{L^p(\sigma B)}2^{kp/p^*},
 \qquad k\geq k_o,
\end{equation}
and, when $1\leq p<s$, the sharper estimate
\begin{equation}
\label{eq29}
 \sup_{E_k\cap B}|u-\gamma|
 \leq C_s\left(
 \frac RL+\left(\frac{s-1}{s-p}\right)^{1/s'}
 \right)b^{-1/s}\mu(B)^{-1/p^*}
 \Vert g\Vert_{L^p(\sigma B)}2^{kp/p^*}.
\end{equation}


To estimate the second term on the right-hand side of \eqref{eq27},
we argue as in the proofs of \eqref{eq10} and \eqref{eq15} when
$s>1$, using H\"older's inequality followed by a geometric series
estimate. When $0<s\leq1$, we use the argument in \eqref{eq39}.
In either case, we obtain
\begin{equation}
\label{eq:local-subcritical-increments}
\begin{aligned}
&2^{1+1/s}b^{-1/s}\mu(B)^{-1/p}
 \Vert g\Vert_{L^p(\sigma B)}
 \sum_{j=k_o+1}^k2^jm_{j-1}^{1/s}\\
&\qquad\leq
 C_{s,p}b^{-1/s}\mu(B)^{-1/p^*}
 \Vert g\Vert_{L^p(\sigma B)}2^{kp/p^*}.
\end{aligned}
\end{equation}
As in the proof of \eqref{eq15}, for $1\leq p<s$, the constant $C_{s,p}$ in this bound
satisfies
\begin{equation}
\label{eq:local-subcritical-increments - constant}
 C_{s,p}\leq
 32\left(\frac{s-1}{s-p}\right)^{1/s'}.
\end{equation}

It remains to estimate the term $a_{k_o}$ in \eqref{eq27}.
Note that, in the proof of Theorem~\ref{T1}, $k_o$ depends only on $p$,
whereas here its choice also involves the ratio $\mu(B)/(bL^s)$.
To express the bound for $a_{k_o}$ in \eqref{eq25} in a form
that can be combined with \eqref{eq:local-subcritical-increments},
we use the lower bound in \eqref{eq24}, which gives
\[
 b^{-1/s}\mu(B)^{-1/p^*}2^{k_op/p^*}
 \geq \Lambda_{s,p}^{p/p^*}b^{-1/p}L^{1-s/p}.
\]
Combining this with \eqref{eq25}, and using $k\geq k_o$ and 
$1-p/p^*=p/s$%$\Lambda_{s,p}>1$
, yields
\[
 a_{k_o}\leq 16\Lambda_{s,p}^{p/s}\frac RL\,b^{-1/s}\mu(B)^{-1/p^*}
 \Vert g\Vert_{L^p(\sigma B)}2^{kp/p^*}.
\]
This proves \eqref{eq28}. 
For $1\leq p<s$, we have
\[
 \Lambda_{s,p}^{p/s}
 =\frac{2^{1/s}}{1-2^{-p/s}}
 \leq\frac{2^{1/s}}{1-2^{-1/s}}.
\]
Consequently, the estimate for $a_{k_o}$ holds with some constant $C_s$
that only depends on $s$.
Combining this estimate with
\eqref{eq:local-subcritical-increments} and
\eqref{eq:local-subcritical-increments - constant}
proves \eqref{eq29}.



We now complete the proof of the subcritical embedding. Set
\begin{equation}
\label{eq:F sets}
F_{k_o}=B\cap E_{k_o},
 \qquad
 F_k=B\cap(E_k\setminus E_{k-1}),\quad k>k_o.
\end{equation}
Up to a null set, we have
\[
B=\bigcup_{k\geq k_o}F_k.
\]
 As in the proof of
\eqref{eq11},
\[
 \sum_{k>k_o}2^{kp}\mu(F_k)
 \leq 2^p\mu(B)\Vert g\Vert_{L^p(\sigma B)}^{-p}
       \int_B g^p\,d\mu
 \leq 2^p\mu(B).
\]
For the $k_o$ level, \eqref{eq24} gives instead
\[
 2^{k_op}\mu(F_{k_o})
 \leq 2^p\Lambda_{s,p}^p\frac{\mu(B)^2}{bL^s}.
\]
Together with \eqref{eq23}, these estimates imply
%
\begin{equation}
\label{eq30}
 \sum_{k\geq k_o}2^{kp}\mu(F_k)
 \leq 2^{p+1}\Lambda_{s,p}^p\frac{\mu(B)^2}{bL^s}.
\end{equation}
As before, when $1\leq p<s$, we can bound the constant $2^{p+1}\Lambda_{s,p}^p$ by some $C_s$.


Summing \eqref{eq28} to the power $p^*$ over the sets $F_k$, exactly as
after \eqref{eq11}, and then using \eqref{eq30}, gives
\[
 \fint_B |u-\gamma|^{p^*}\,d\mu
 \leq C_{s,p}^{p^*}2^{p+1}\Lambda_{s,p}^p\left(\frac RL\right)^{p^*}
 b^{-p^*/s-1}L^{-s}\Vert g\Vert_{L^p(\sigma B)}^{p^*},
\]
where $C_{s,p}$ is the constant in \eqref{eq28}.
Since $1/s+1/p^*=1/p$ and $-s/p^*=1-s/p$, it follows that
\[
 \left(\fint_B|u-\gamma|^{p^*}\,d\mu\right)^{1/p^*}
 \leq C_{s,p}\left(2^{p+1}\Lambda_{s,p}^p\right)^{1/p^*}\frac RL\,b^{-1/p}L^{1-s/p}
 \Vert g\Vert_{L^p(\sigma B)}.
\]
This proves the first assertion in the subcritical case.

For $1\leq p<s$, we repeat this calculation using
\eqref{eq29} in place of \eqref{eq28}.
The factor coming from \eqref{eq30} satisfies
\[
 \left(2^{p+1}\Lambda_{s,p}^p\right)^{1/p^*}
 \leq
 \left(\frac{2^{s+2}}{(1-2^{-1/s})^s}\right)^{1/p^*}
 \leq
 \left(\frac{2^{s+2}}{(1-2^{-1/s})^s}\right)^{1/s'},
\]
where the last inequality uses $p^*\geq s'$.
Thus this factor is bounded independently of $p$ in the
range $1\leq p<s$, and we obtain
\[
 \left(\fint_B|u-\gamma|^{p^*}\,d\mu\right)^{1/p^*}
 \leq C_s\left(
 \frac RL+\left(\frac{s-1}{s-p}\right)^{1/s'}
 \right)b^{-1/p}L^{1-s/p}
 \Vert g\Vert_{L^p(\sigma B)}.
\]
Finally, whenever $p^*\geq1$, the same averaging argument
as in Theorem~\ref{T1} replaces $\gamma$ by $u_B$ at the
cost of a factor of $2$. This completes the proof of the theorem
in the subcritical case.


\medskip
\noindent
{\sc The critical case $p=s>1$.}
Write $s'=s/(s-1)$. Repeating the H\"older estimate leading to
\eqref{eq12}, now starting at the local level $k_o$ instead of
at $1$, and using \eqref{eq22}, \eqref{eq25}, and \eqref{eq27},
we obtain
\begin{equation}
\label{eq31}
 \sup_{E_k\cap B}|u-\gamma|
 \leq C_sb^{-1/s}\Vert g\Vert_{L^s(\sigma B)}
 \left(\frac RL+(k-k_o)^{1/s'}\right),
 \qquad k>k_o.
\end{equation}
The same estimate, with $k=k_o$, follows from \eqref{eq25}.

With the sets $F_k$ defined as in \eqref{eq:F sets}, observe that
since $k_o\geq1$, \eqref{eq21} gives
\begin{equation}
\label{eq32}
 \frac{\mu(F_k)}{\mu(B)}
 \leq2^{-s(k-1)}
 \leq2^{-s(k-k_o)},
 \qquad k>k_o.
\end{equation}
By \eqref{eq31} and $L/R\leq1$, we have
\[
 \left(
 \frac LR\frac{b^{1/s}|u-\gamma|}
 {\Vert g\Vert_{L^s(\sigma B)}}
 \right)^{s'}
 \leq C_s(1+k-k_o)
 \qquad\text{on }F_k,\quad k\geq k_o.
\]
As in the argument following \eqref{eq12}, choosing $c_s>0$
sufficiently small and summing over the sets $F_k$ yields
%
\begin{equation}
\label{eq33}
 \fint_B
 \exp\!\left[
 c_s\left(
 \frac LR\frac{b^{1/s}|u-\gamma|}
 {\Vert g\Vert_{L^s(\sigma B)}}
 \right)^{s'}
 \right]d\mu
 \leq C_s.
\end{equation}
Indeed, for sufficiently small $c_s$, the integrand in
\eqref{eq33} is at most $C_s2^{s(k-k_o)/2}$ on $F_k$,
$k>k_o$, and is bounded by $C_s$ on $F_{k_o}$ by
\eqref{eq25}. Together with \eqref{eq32}, these bounds
show that the integral average in \eqref{eq33} is at most
\[
 C_s\left(
 1+\sum_{k=k_o+1}^{\infty}
 2^{-s(k-k_o)}2^{s(k-k_o)/2}
 \right)
 =\frac{C_s}{1-2^{-s/2}}.
\]
In particular, $u\in L^1(B)$. The same argument used in the proof of Lemma~\ref{T6} with $X$ replaced by $B$ allows us to replace $\gamma$ by $u_B$ in \eqref{eq33}, after replacing
$c_s$ by $c_s/2^{s'}$. This proves the theorem in the case
$p=s>1$.


\medskip
\noindent
{\sc A pointwise estimate for the remaining cases.}
We first establish a pointwise estimate that will be used to prove uniform
continuity when $p=s\leq1$ and when $p>s$, as well as the H\"older
estimate when $p>s$.

Fix $k\geq k_o$ and $x\in B\setminus N$. 
We now choose a point $x_k\in E_k$ as follows:
If $x\in E_k$, put $x_k=x$.
Otherwise, since $\bigcup_{m\geq k_o}E_m=\sigma B\setminus N$,
we may choose the smallest $m>k$ such that $x\in E_m$. Starting with $x_m=x$,
apply the chain construction down to level $k$, choosing
\[
 x_{j-1}\in B(x_j,r_j)\cap E_{j-1},
 \qquad j=m,m-1,\ldots,k+1,
\]
which gives us a point $x_k\in E_k$. Note that the estimates along the chain give
\[
 d(x,x_k)\leq\sum_{j=k+1}^m r_j
\]
and
\[
 |u(x)-u(x_k)|
 \leq 2\mu(B)^{-1/p}\Vert g\Vert_{L^p(\sigma B)}
       \sum_{j=k+1}^m2^jr_j.
\]
For $x,y\in B\setminus N$, choose $x_k,y_k\in E_k$ in this way.
Since $u|_{E_k}$ is Lipschitz with constant at most
$2^{k+1}\mu(B)^{-1/p}\Vert g\Vert_{L^p(\sigma B)}$, the triangle
inequality gives
\[
\begin{split}
 |u(x)-u(y)|
 &\leq |u(x)-u(x_k)|+|u(y)-u(y_k)|\\
 &\quad+2^{k+1}\mu(B)^{-1/p}\Vert g\Vert_{L^p(\sigma B)}
 \bigl(d(x,x_k)+d(x,y)+d(y,y_k)\bigr).
\end{split}
\]
Using the preceding bounds, together with
$2^k\sum_{j=k+1}^{\infty}r_j\leq\sum_{j=k+1}^{\infty}2^jr_j$, we obtain
\begin{equation}
\label{eq:T2-two-point}
\begin{split}
 |u(x)-u(y)|
 &\leq C_s b^{-1/s}\mu(B)^{-1/p}\Vert g\Vert_{L^p(\sigma B)}
       \sum_{j=k+1}^{\infty}2^jm_{j-1}^{1/s}\\
 &\quad+2^{k+1}\mu(B)^{-1/p}\Vert g\Vert_{L^p(\sigma B)}d(x,y).
\end{split}
\end{equation}
If the series tail in \eqref{eq:T2-two-point} tends to zero as
$k\to\infty$, then $u|_{B\setminus N}$ is uniformly continuous.
Indeed, if this is the case then, for each $\varepsilon>0$, we can choose $k$ so that the first term is less than $\varepsilon/2$, and then take $d(x,y)$ small enough so that the second term is less than $\varepsilon/2$. Since $B\setminus N$ is dense in $B$, this gives a unique uniformly continuous representative on $B$, which we continue to denote by $u$, satisfying \eqref{eq:T2-two-point} for all $x,y\in B$.


\medskip
\noindent
{\sc The critical case $p=s$, $0<s\leq1$.}
As in the corresponding case of Theorem~\ref{T1}, we use
$\ell^s\subset\ell^1$. Together with \eqref{eq22}, this gives
\begin{equation}
\label{eq:T2-critical-tail}
 \sum_{j=k+1}^{\infty}2^jm_{j-1}^{1/s}
 \leq
 \left(\sum_{j=k+1}^{\infty}2^{js}m_{j-1}\right)^{1/s}
 \longrightarrow0
 \qquad\text{as }k\to\infty.
\end{equation}
Thus the preceding argument gives a uniformly continuous
representative on $B$.
Moreover, \eqref{eq22} bounds the right-hand side of
\eqref{eq:T2-critical-tail} at $k=k_o$ by $C_s\mu(B)^{1/s}$.
Combining this bound with \eqref{eq25} and \eqref{eq27}, and using
$R/L\geq1$, yields
\begin{equation}
\label{eq34}
 \sup_{E_k\cap B}|u-\gamma|
 \leq C_s\frac RL\,b^{-1/s}\Vert g\Vert_{L^s(\sigma B)},
 \qquad k\geq k_o.
\end{equation}
Since $B\setminus N=\bigcup_{k\geq k_o}(E_k\cap B)$,
continuity extends this bound to all of $B$. Consequently,
\[
 \operatorname{osc}_B u
 \leq 2\sup_B|u-\gamma|
 \leq C_s\frac RL\,b^{-1/s}\Vert g\Vert_{L^s(\sigma B)}.
\]
This completes the critical case.


\medskip
\noindent
{\sc The supercritical case $p>s$.}
By \eqref{eq21}, for $k\geq k_o$,
\begin{equation}
\label{eq37}
\begin{split}
 \sum_{j=k+1}^{\infty}2^jm_{j-1}^{1/s}
 &\leq 2^{p/s}\mu(B)^{1/s}
        \sum_{j=k+1}^{\infty}2^{j(1-p/s)}\\
 &\leq C_{s,p}\mu(B)^{1/s}2^{k(1-p/s)}.
\end{split}
\end{equation}
Since $1-p/s<0$, these tails tend to zero. Hence the argument
following \eqref{eq:T2-two-point} again gives a continuous
representative of $u$ on $B$, which we continue to denote by $u$.


As in the supercritical case of Theorem~\ref{T1}, we now sum the
increments. Equations \eqref{eq27} and \eqref{eq37} give
\[
 |u(x)-\gamma|\leq a_{k_o}
 +C_{s,p}b^{-1/s}\mu(B)^{1/s-1/p}
 \Vert g\Vert_{L^p(\sigma B)}2^{k_o(1-p/s)},
 \qquad x\in B\setminus N.
\]
The initial term is bounded by \eqref{eq25}. Since $1-p/s<0$,
the lower bound in \eqref{eq24} yields
\[
 b^{-1/s}\mu(B)^{1/s-1/p}2^{k_o(1-p/s)}
 \leq \Lambda_{s,p}^{1-p/s}b^{-1/p}L^{1-s/p}.
\]
Using $R/L\geq1$ and extending by continuity, we obtain
\begin{equation}
\label{eq38}
 |u-\gamma|
 \leq C_{s,p}\frac RL\,b^{-1/p}L^{1-s/p}
 \Vert g\Vert_{L^p(\sigma B)}
 \qquad\text{on }B.
\end{equation}
The stated oscillation estimate now follows from
$\operatorname{osc}_B u\leq2\sup_B|u-\gamma|$.

We now prove the H\"older estimate.
Let $x,y\in B$ be distinct, and set
$c_{s,p}=(2\Lambda_{s,p})^{-p/s}$.
If $d(x,y)\leq c_{s,p}L$, \eqref{eq24} implies that there exists
$k\geq k_o$ such that
%
\begin{equation}
\label{eq:T2-holder-level}
 2^k\leq b^{-1/p}\mu(B)^{1/p}d(x,y)^{-s/p}<2^{k+1}.
\end{equation}
%
Using \eqref{eq37} in \eqref{eq:T2-two-point}, we obtain
%
\begin{equation}
\label{eq:T2-holder-two-terms}
\begin{split}
 |u(x)-u(y)|
 &\leq C_{s,p}b^{-1/s}\mu(B)^{1/s-1/p}
       \Vert g\Vert_{L^p(\sigma B)}2^{k(1-p/s)}\\
 &\quad+2^{k+1}\mu(B)^{-1/p}
       \Vert g\Vert_{L^p(\sigma B)}d(x,y).
\end{split}
\end{equation}
The right-hand inequality in \eqref{eq:T2-holder-level} bounds the
first term in \eqref{eq:T2-holder-two-terms}, since $1-p/s<0$,
while the left-hand inequality bounds the second. Consequently,
\[
 |u(x)-u(y)|
 \leq C_{s,p}b^{-1/p}d(x,y)^{1-s/p}
       \Vert g\Vert_{L^p(\sigma B)}.
\]
If $d(x,y)>c_{s,p}L$, \eqref{eq38} immediately gives
\[
\begin{split}
 |u(x)-u(y)|
 &\leq C_{s,p}\frac RL\,b^{-1/p}L^{1-s/p}
       \Vert g\Vert_{L^p(\sigma B)}\\
 &\leq C_{s,p}\frac RL\,b^{-1/p}d(x,y)^{1-s/p}
       \Vert g\Vert_{L^p(\sigma B)}.
\end{split}
\]
Combining these last two estimates and using $R/L\geq1$ proves the desired
H\"older estimate, and the proof of Theorem~\ref{T2} is complete.
\end{proof}

\end{document}
